\paragraph{Recovery from a single high band.}
There is a stronger reconstruction which discards every zero of height
below $T$. Take a real even $w\in C_c^\infty(\{s:1<|s|<2\})$ with
$\int_{\mathbb R} w(s)\,ds=2\pi$. For every fixed $u\in\mathbb R$
and every $A>0$,
\begin{equation}
 J(u)=-\frac1T\sum_{T<|\gamma|<2T}
          m(\gamma) w(\gamma/T)e^{i\gamma u}
       +O_{q,d,u,w,A}(T^{-A}).
 \label{eq:packet-band}
\end{equation}
Thus a sequence of disjoint high bands gives, at $u=0$, successively
accurate readings of $\log d$. At each fixed prime it gives readings of
the corresponding prime jump. The error statement concerns exact signed
multiplicities and exact ordinates.

To prove \eqref{eq:packet-band}, take
$g(v)=(2\pi)^{-1}\int w(s)e^{isv}\,ds$. This Schwartz function satisfies
$g(0)=1$ and $\int v^j g(v)\,dv=0$ for every integer $j\ge0$, since $w$
vanishes near zero. Pair $\mu=F''-F/4$ against $g(T(t-u))$.
Its Fourier series as a distribution gives exactly the negative finite
sum on the right of \eqref{eq:packet-band}. To evaluate the same pairing
locally, multiply by a smooth cutoff equal to one near $u$ with support
containing no other atom. The atom contributes $J(u)$. The smooth compactly
supported density contributes $O_A(T^{-A})$ by Taylor expansion and the
vanishing moments. For the complementary cutoff, integrate twice by parts
onto the test function using $\lVert F\rVert_\infty<\infty$.
Away from $u$, the $L^1$ norms of this function and its first two derivatives
are $O_A(T^{-A})$ by Schwartz decay. This last step permits localization
even though the positive prime masses and the smooth density on the
right of \eqref{eq:packet-ode} individually grow exponentially at infinity.

This formula also quantifies errors in the spectral data. Suppose the signed
entries contributing near the band have ordinate errors at most
$\epsilon_T\le1$. Allow deletions, insertions and weight errors with total
absolute weight $E_T$ in the enlarged band $T/2<|\gamma|<3T$.
In the estimator of \eqref{eq:packet-band} the resulting error is at most
\[
 O_{q,d,w}\!\left(\frac{E_T}{T}
       +\epsilon_T (|u|+1/T)\log(2+T)\right).
\]
Indeed $w(v/T)e^{ivu}$ is Lipschitz with constant
$\lVert w'\rVert_\infty/T+|u|\lVert w\rVert_\infty$, while the sum of
absolute weights in the enlarged band is $O_{q,d}(T\log(2+T))$ by the
zero count. The cutoff has compact support, so a displacement of at most
one can only bring entries from that enlarged band into the estimator.
Thus $E_T=o(T)$ and $\epsilon_T=o(1/\log T)$ preserve every fixed jump
in the limit. This uses signed entries and does not require pairing two
nearby ordinates before summation. In particular, two distinct parameter
pairs have absolute signed multiplicity discrepancy $\gg T$ in each
sufficiently high dyadic band: their jump functions differ somewhere, and
\eqref{eq:packet-band} at that point bounds the discrepancy from below.
